% Written in AMS-LaTeX   version 4/29/99



\documentstyle[12pt,twoside]{amsart}
%\documentstyle[amstex,12pt,leqno]{book}
%\documentstyle[amstex,12pt,leqno]{amsart}
%\documentstyle[12pt,leqno]{amsbook}
%\title{Title of the paper}
%\author{Author} 
%\date{Date}



        



\let \cedilla =\c
\renewcommand{\c}[0]{{\mathbb C}}  
\let \crossedo =\o
\renewcommand{\o}[0]{{\cal O}} 
\newcommand{\z}[0]{{\mathbb Z}}
\newcommand{\n}[0]{{\mathbb N}}
\let \ringaccent=\r  %%% \r  shorthand for 'ring accent'
\renewcommand{\r}[0]{{\mathbb R}} 

\renewcommand{\a}[0]{{\mathbb A}} 


\newcommand{\p}[0]{{\mathbb P}}
\newcommand{\f}[0]{{\mathbb F}}
\newcommand{\q}[0]{{\mathbb Q}}
\newcommand{\map}[0]{\dasharrow}
\newcommand{\qtq}[1]{\quad\mbox{#1}\quad}
\newcommand{\spec}[0]{\operatorname{Spec}}
\newcommand{\pic}[0]{\operatorname{Pic}}
\newcommand{\gal}[0]{\operatorname{Gal}}
\newcommand{\cont}[0]{\operatorname{cont}}
\newcommand{\rank}[0]{\operatorname{rank}}
\newcommand{\mult}[0]{\operatorname{mult}}
\newcommand{\discrep}[0]{\operatorname{discrep}}
\newcommand{\totaldiscrep}[0]{\operatorname{totaldiscrep}}

\newcommand{\supp}[0]{\operatorname{Supp}}    
\newcommand{\red}[0]{\operatorname{red}}    
\newcommand{\codim}[0]{\operatorname{codim}}    
\newcommand{\im}[0]{\operatorname{im}}    
\newcommand{\proj}[0]{\operatorname{Proj}}    
\newcommand{\wt}[0]{\operatorname{wt}}    
\newcommand{\socle}[0]{\operatorname{socle}}    
\newcommand{\coker}[0]{\operatorname{coker}}    
\newcommand{\ext}[0]{\operatorname{Ext}}    
\newcommand{\Hom}[0]{\operatorname{Hom}}    
\newcommand{\trace}[0]{\operatorname{Trace}}  
\newcommand{\cent}[0]{\operatorname{center}}
\newcommand{\bs}[0]{\operatorname{Bs}}
\newcommand{\aut}[0]{\operatorname{Aut}}    
\newcommand{\specan}[0]{\operatorname{Specan}}    
\newcommand{\inter}[0]{\operatorname{Int}}    
\newcommand{\sing}[0]{\operatorname{Sing}}    
\newcommand{\ex}[0]{\operatorname{Ex}}    
\newcommand{\chr}[0]{\operatorname{char}}    




\newcommand{\nec}[1]{\overline{NE}({#1})}

\newcommand{\rup}[1]{\lceil{#1}\rceil}
\newcommand{\rdown}[1]{\lfloor{#1}\rfloor}

\newsymbol\subsetneq 2328
\newsymbol\onto 1310
\def\into{\DOTSB\lhook\joinrel\rightarrow}
\newsymbol\twoheadrightarrow 1310


% Definitions for new environments
% theorem style plain --- default
%\newtheorem{thm}{Theorem}[chapter]
\newtheorem{thm}{Theorem}%[section]
\newtheorem{mainthm}[thm]{Main Theorem}
\newtheorem{question}[thm]{Question}
\newtheorem{lem}[thm]{Lemma}
\newtheorem{cor}[thm]{Corollary}
\newtheorem{cors}[thm]{Corollaries}
\newtheorem{prop}[thm]{Proposition}
\newtheorem{crit}[thm]{Criterion}
\newtheorem{conj}[thm]{Conjecture}
\newtheorem{principle}[thm]{Principle} %!!!!!!!!!!!!!!!!!!!!!!
\newtheorem{complement}[thm]{Complement}%!!!!!!!!!!!!!!!!!!!!!!

\theoremstyle{definition}
\newtheorem{defn}[thm]{Definition}
\newtheorem{condition}[thm]{Condition}
\newtheorem{say}[thm]{}
\newtheorem{exmp}[thm]{Example}
\newtheorem{hint}[thm]{Hint}
\newtheorem{exrc}[thm]{Exercise}
\newtheorem{prob}[thm]{Problem}
\newtheorem{const}[thm]{Construction}   %!!!!!!!!!!!!!!!!
\newtheorem{ques}[thm]{Question}    %!!!!!!!!!!!!!!!!!!!!
\newtheorem{alg}[thm]{Algorithm}
\newtheorem{rem}[thm]{Remark}          
%\renewcommand{\theremark}{} 
\newtheorem{note}[thm]{Note}            %\renewcommand{\thenote}{} 
\newtheorem{summ}[thm]{Summary}         %\renewcommand{\thesumm}{} 
\newtheorem{ack}{Acknowledgments}       \renewcommand{\theack}{} 
\newtheorem{notation}[thm]{Notation}   
\newtheorem{defn-lem}[thm]{Definition--Lemma}




\theoremstyle{remark}
\newtheorem{claim}[thm]{Claim}  
\newtheorem{case}{Case} 
\newtheorem{subcase}{Subcase}   
\newtheorem{step}{Step}
\newtheorem{approach}{Approach}
%\newtheorem{principle}{Principle}
\newtheorem{fact}{Fact}
\newtheorem{subsay}{}

%%%%%%%%%%%%%%%%%%%%%%%%%%%%%%%%%%%%%%%%%%%%%%%%%%%%%%%%%%%%%
\setcounter{section} {0}
 %\documentstyle[12pt]{book}
%\setlength{\topmargin}{0in}
%\setlength{\headsep}{0in}
%\setlength{\headheight}{0in}
%\setlength{\textheight}{8.25in}
%\setlength{\parindent}{0in}
%\pagestyle{empty}


\begin{document}
\bibliographystyle{amsplain}



\title{Fixed points of group actions and rational maps}
\author{J{\'a}nos Koll{\'a}r and Endre Szab{\'o}}


\maketitle


The aim of this note is to give simple proofs of
the results in section 5 of
\cite{ry} about the behaviour of fixed points of finite group
actions under rational maps. Our proofs   work in any
characteristic.


\begin{lem}\label{class.defn}
 Let $K$ be an algebraically closed field and  $H$
a (not necessarily connected) linear
algebraic group over $K$.  The following are equivalent.
\begin{enumerate}
\item Every representation  $H\to GL(n,K)$ has an
$H$-eigenvector.
\item There is a  (not necessarily connected) unipotent, normal 
subgroup $U<H$ such that $H/U$ is abelian.
\end{enumerate}
\end{lem}

Proof. Let $H\to GL(n,K)$ be a faithful representation.
If (\ref{class.defn}.1) holds then $H$ is conjugate to
an upper triangular subgroup, this implies (\ref{class.defn}.2).

Conversely, any representation of a unipotent group
has eigenvectors (with eigenvalue 1, cf.\ \cite[I.4.8]{borel})
and the subspace of all eigenvectors is an
$H/U$-representation.\qed

\begin{prop}[Going down]\label{Going-down}
  Let $K$ be an algebraically closed
field,
$H$ a linear  algebraic group over $K$ and $f:X\map Y$ an
$H$-equivariant map of
$K$-schemes. Assume that
\begin{enumerate}
\item $H$ satisfies the equivalent conditions
of (\ref{class.defn}),
\item $H$ has a smooth fixed point on $X$, and
\item $Y$ is proper.
\end{enumerate}
\noindent Then $H$ has a fixed point on $Y$.
\end{prop}

Proof. The proof is by induction on $\dim X$. The  case  $\dim
X=0$ is clear.

Let $x\in X$ be a smooth $H$-fixed point and consider the blow up
$B_xX$ with exceptional divisor $E\cong \p^{n-1}$. 
The $H$-action lifts to $B_xX$ and so we get an $H$-action on $E$
which has a fixed point by (\ref{class.defn}.1).  Since $Y$ is
proper, the induced rational map $B_xX\to X\map Y$ is defined
outside a subset of codimension at least 2. Thus we get an
$H$-equivariant rational map $E\map Y$. By induction, there is a 
fixed point on $Y$.\qed


\begin{rem} If $H$ does not satisfy the  conditions
of (\ref{class.defn})
then (\ref{Going-down}) fails for some actions.
Indeed, let $H\to GL(n,K)$ be a representation without an 
$H$-eigenvector.  This gives an $H$-action on $\p^n$ with a single fixed
point $Q\in \p^n$. The corresponding action on $B_Q\p^n$ has no fixed
points.
\end{rem}

\begin{prop}[Going up]\label{Going-up}
  Let $K$ be an algebraically closed field
and
$H$ a finite abelian group of prime power order $q^n$ 
($q$ is allowed to coincide with $\chr K$). Let $p:X\map Z$ be an
$H$-equivariant map of irreducible
$K$-schemes. Assume that
\begin{enumerate}
\item $f$ is generically finite, dominant and $q\not\vert\deg (X/Z)$,
\item $H$ has a smooth fixed point on $Z$, and
\item $X$ is proper.
\end{enumerate}
\noindent Then $H$ has a fixed point on $X$.
Moreover, if $X\map Y$ is an $H$-equivariant
map to a proper $K$-scheme then $H$ has a fixed point on $Y$.
\end{prop}

Proof.  The proof is by induction on $\dim Z$. 
The  case  $\dim Z=0$ is clear.

Let $z\in Z$ be a smooth fixed point and $E\subset B_zZ$
the exceptional divisor. Let $\bar p:\bar X\to B_zZ$
denote the normalization of $B_zZ$ in the field of rational
functions of $X$ and $F_i\subset \bar X$ the divisors lying over
$E$.  $H$ acts on the set   $\{F_i\}$. Let ${\cal F}_j$
denote the   $H$-orbits and in each pick a divisor
$F^*_j\in {\cal F}_j$.
By the ramification formula,
$$
\deg (X/Z)=\sum_j 
|{\cal F}_j|\cdot\deg (F^*_j/E)\cdot e(\bar p, F^*_j)
$$
 where
$e(\bar p,F^*_j)$ denotes
the ramification index of $\bar p$ at the generic point of $F^*_j$.
 Since $\deg (X/Z)$ is not divisible by $q$,
there is an orbit ${\cal F}_0$ consisting of  a single element
$F^*_0$  such that $\deg (F^*_0/E)$ is not divisible by $q$.

We have $H$-equivariant rational maps 
$F^*_0\map E$, $F^*_0\map Z$ and $F^*_0\map Y$.
By induction $H$ has a fixed point on $Y$, and also on $Z$.
\qed

\begin{rem} We see from the proof that (\ref{Going-up})
also holds  if $H$ is abelian and only one of the prime divisors of
$|H|$ is less than $\deg (X/Z)$.
\end{rem}


The method   also gives a simpler proof of a result
of \cite{nishi}. One can view this as a version of
(\ref{Going-down}) where $H$ is the absolute Galois group of $K$.


\begin{prop}[Nishimura lemma]\label{nishimura}
  Let $K$ be a 
field  and $f:X\map Y$ a 
map of
$K$-schemes. Assume that
\begin{enumerate}
\item $X$ has a smooth $K$-point, and
\item $Y$ is proper.
\end{enumerate}
\noindent Then $Y$ has a   $K$-point.
\end{prop}

Proof. The proof is by induction on $\dim X$. The  case  $\dim
X=0$ is clear.

Let $x\in X$ be a smooth $K$-point and consider the blow up
$B_xX$ with exceptional divisor $E\cong \p^{n-1}$. 
$E$ has smooth $K$-points.
 Since $Y$ is
proper, the induced rational map $B_xX\to X\map Y$ is defined
outside a subset of codimension at least 2 and we get a rational
map $E\map Y$. By induction, there is a  $K$-point on $Y$.\qed


\begin{rem}  One can combine (\ref{Going-down}) and 
(\ref{nishimura}) if we know that any $H$-representation has an
eigenvector defined over $K$. There are two interesting cases
where this condition holds:
\begin{enumerate}
\item  $H$ is Abelian of order $n$ and $K$ contains all $n$th
roots of unity.
\item $H$ is nilpotent and its order is a power of $\chr K$.
\end{enumerate}
\end{rem}


\begin{thebibliography}{Chislenko88}

\bibitem[Borel91]{borel}  A.\ Borel, Linear algebraic groups,
second ed. Springer, 1991


\bibitem[Nishimura55]{nishi} H.\ Nishimura, Some remarks on
rational points, Mem. Coll. Sci. Univ. Kyoto  29 (1955)
189-192 


\bibitem[Reichstein-Youssin99]{ry} Z.\ Reichstein and B.\
Youssin, Essential dimensions of algebraic groups and a resolution
theorem for $G$-varieties (math.AG/9903162)


\end{thebibliography}

\vskip1cm

\noindent University of Utah, Salt Lake City UT 84112 

\begin{verbatim}kollar@math.utah.edu\end{verbatim}


\noindent Mathematical Institut, Budapest, PO.Box 127, 1364 Hungary

\begin{verbatim}endre@math-inst.hu\end{verbatim}



\end{document}


